% !TEX root = ../main.tex
%==============================================================
%  ABSTRACT (English)
%==============================================================
\thispagestyle{empty}
\begin{center}
{\Large\bfseries ABSTRACT}
\end{center}
\vspace{1cm}

Computer-generated images play an important role in many real-world applications,
ranging from recreational fields such as cinema and video games to engineering
fields such as architecture and medicine. Physically based rendering is the class
of algorithms that codifies the physical laws of light transport into mathematical
models in order to produce highly realistic images. Among these, Path Tracing is a
general and accurate method, but it is computationally expensive because it must
simulate the paths of millions of light rays through Monte Carlo integration,
resulting in long rendering times and making interactive frame rates difficult to
achieve.

In this project, I study and build a Path Tracer that runs in parallel on NVIDIA's
CUDA architecture to exploit the computational power of the GPU and achieve
real-time rendering. The system implements the complete ray tracing pipeline on the
GPU, including: the Path Tracing algorithm with temporal accumulation and
anti-aliasing; a material system supporting diffuse, metallic, and dielectric
surfaces; a thin-lens camera model with depth-of-field and motion-blur effects; a
CUDA--OpenGL interoperability mechanism that enables display directly from video
memory; together with a Mitsuba~3 scene loader and an interactive control interface.

Experimental results show that the system is able to render the test scenes at
interactive rates while progressively converging to a low-noise image when the
camera remains static. The project also identifies future optimization and
extension directions, such as bounding volume hierarchy (BVH) acceleration
structures and advanced denoising techniques.

\vspace{1.5cm}
\begin{flushright}
Student\\[0.3cm]
\textit{(Signature and full name)}
\end{flushright}
